\documentclass{article}
\usepackage{graphicx}
\usepackage{booktabs}
\usepackage{amsmath}
\usepackage{hyperref}

\title{Failure Propagation in Agentic RAG Pipelines: Do Retrieval Errors Compound, Correct, or Isolate?}
\author{[Intern Name], Spurthi Setty \\
Anote AI}
\date{}

\begin{document}

\maketitle

\begin{abstract}
\noindent
Retrieval-Augmented Generation (RAG) pipelines are increasingly embedded within multi-step agentic workflows, yet the downstream effects of retrieval errors on pipeline-level outcomes remain unstudied. We conduct a systematic failure injection study on multi-step agentic RAG pipelines, classifying retrieval error propagation into three patterns: \textit{compounding} (error amplifies at each step), \textit{self-correcting} (downstream steps recover), and \textit{isolated} (error does not affect subsequent steps). Using a controlled testbed with five pipeline configurations and three domains (finance, legal, medical), we find that X\% of retrieval errors compound, Y\% self-correct, and Z\% remain isolated. We identify the pipeline stages most vulnerable to error propagation and propose two intervention mechanisms that reduce compounding by A\%.
\end{abstract}

\section{Introduction}
% RAG is widely deployed; single-step analysis is insufficient for agentic pipelines
% Error propagation in multi-step systems is a known risk (cite)
% This paper: empirical taxonomy of failure propagation in agentic RAG

\section{Related Work}
\subsection{RAG Systems and Failure Modes}
% Cite arXiv:2404.07221 (our prior work)
\subsection{Multi-Step Agentic Systems}
\subsection{Error Propagation in NLP Pipelines}

\section{Agentic RAG Pipeline Testbed}
\subsection{Pipeline Architecture}
% 5 pipeline configurations: 2-step, 3-step, 4-step with/without re-ranking
\subsection{Failure Injection Methodology}
% Controlled retrieval error types: wrong chunk, partial chunk, hallucinated chunk, empty retrieval
\subsection{Propagation Tracking}
% How we trace errors through pipeline steps

\section{Failure Propagation Taxonomy}
\subsection{Compounding Errors}
\subsection{Self-Correcting Errors}
\subsection{Isolated Errors}

\section{Experimental Setup}
\subsection{Domains}
% Finance (FinanceBench), Legal (CUAD), Medical (PubMedQA)
\subsection{Evaluation Metrics}
% Final answer accuracy, intermediate step accuracy, propagation rate

\section{Results}
% Figure: Error propagation distribution by pipeline configuration
% Table: Propagation rates by error type and domain
% Figure: Intervention effectiveness

\subsection{Propagation Rates by Error Type}
\subsection{Domain Comparison}
\subsection{Intervention Effectiveness}

\section{Analysis}
\subsection{Most Vulnerable Pipeline Stages}
\subsection{Why Some Errors Self-Correct}

\section{Conclusion}

\section*{References}
\small
\begin{enumerate}
    \item Setty et al. Improving Retrieval for RAG based Question Answering Models on Financial Documents, 2024. arXiv:2404.07221.
    \item Lewis et al. Retrieval-Augmented Generation for Knowledge-Intensive NLP Tasks, 2020. NeurIPS 2020.
    % Add more references
\end{enumerate}

\end{document}
